\chapter{局部反演}
\label{chap:local-inversion-solutions}

\begin{solution}{neumann-series-inverse-bound}
令 $E=A^{-1}(B-A)$，则 $B=A(I+E)$，且
\[
 B^{-1}=\sum_{k=0}^{\infty}(-E)^kA^{-1}.
\]
取算子范数并使用几何级数，
\[
 \norm{B^{-1}}
 \le\sum_{k=0}^{\infty}\norm E^k\norm{A^{-1}}
 \le\frac{\norm{A^{-1}}}{1-\norm E}.
\]
\end{solution}

\begin{solution}{derivative-close-injective}
设 $G(z)=A^{-1}f(z)-z$。对 $z\in V$，
$DG(z)=A^{-1}Df(z)-I$，故 $\norm{DG(z)}\le q$。连接 $x,y$ 的线段
完全落在凸集 $V$ 中，向量值中值不等式给出
\[
 \norm{G(y)-G(x)}\le q\norm{y-x}.
\]
把 $G(y)-G(x)$ 展开，并分别使用三角不等式及反三角不等式，得到
\[
 (1-q)\norm{y-x}
 \le\norm{A^{-1}(f(y)-f(x))}
 \le(1+q)\norm{y-x}.
\]
凸性只用于保证整条线段都在可使用导数上界的区域内。若区域不凸，
可在其中的小凸邻域内重新取得局部结论；对于原来给定的两个点，
不能假定存在一个既包含它们的连线、又留在定义域内且保持原导数界的小球。
\end{solution}

\begin{solution}{inverse-contraction-ball}
在正规化条件 $f(0)=0$、$Df(0)=I$ 下取
$\sup_{\overline B_r}\norm{Df-I}\le q<1$。于是
\[
 \norm{x-f(x)}\le q\norm x,
\]
并且对 $T_y(x)=x+y-f(x)$，
\[
 \norm{T_y(x)}\le q\norm x+\norm y.
\]
当 $\norm x\le r$、$\norm y\le(1-q)r$ 时，右端不超过 $r$，所以
\[
 T_y(\overline B_r)\subset\overline B_r.
\]
同时 $\norm{T_y(x)-T_y(z)}\le q\norm{x-z}$。因此闭球上的压缩映射定理
保证每个 $y\in B(0,(1-q)r)$ 都有唯一原像。这就是证明直接保证的像球
半径；实际像集可能更大。
\end{solution}

\begin{solution}{inverse-derivative-formula}
若已经知道 $g$ 存在且可微，由 $f(g(y))=y$ 和链式法则得到
\[
 Df(g(y))Dg(y)=I,
\]
从而 $Dg(y)=Df(g(y))^{-1}$。这段计算以 $g$ 已存在、可微为前提，不能
证明邻近每个 $y$ 都有唯一原像，也不能排除一个形式上的解支在邻域内断裂。
逆函数定理中的压缩映射论证正是用来建立这些前提的。
\end{solution}

\begin{solution}{polar-exponential-local-not-global}
Jacobian 为
\[
 Df(x,y)=e^x
 \begin{pmatrix}\cos y&-\sin y\\\sin y&\cos y\end{pmatrix},
 \qquad \det Df=e^{2x}>0.
\]
固定源点 \((x_0,y_0)\)。对目标坐标 \((u,v)\) 作旋转
\[
 U=u\cos y_0+v\sin y_0,\qquad V=-u\sin y_0+v\cos y_0.
\]
在目标点 \((e^{x_0}\cos y_0,e^{x_0}\sin y_0)\) 附近有 \(U>0\)，可取
\[
 x=\log\sqrt{u^2+v^2},\qquad y=y_0+\arctan(V/U)
\]
作为局部逆，源的角度限制在 \(|y-y_0|<\pi/2\)。
原映射满足 \(f(x,y+2\pi)=f(x,y)\)，因此不可能有双侧全局逆。
甚至连续右逆也不存在：若有，其第二分量在单位圆上的复合记为
\(\theta(t)\)，\(0\le t\le2\pi\)。右逆条件意味着
\(\theta(t)-t\in2\pi\Z\)。连续的整数值函数在区间上恒定，故
\(\theta(2\pi)-\theta(0)=2\pi\)。但圆周两端表示同一点，右逆的单值性
又给出 \(\theta(2\pi)=\theta(0)\)，矛盾。
\end{solution}

\begin{solution}{complex-square-local-inverse}
\[
 DF(x,y)=\begin{pmatrix}2x&-2y\\2y&2x\end{pmatrix},
 \qquad\det DF=4(x^2+y^2).
\]
所以除原点外处处局部可逆。原点的每个邻域都含有一对非零相反点，
而 \(F(x,y)=F(-x,-y)\)，故原点不局部可逆。
同一个等式也说明穿孔平面上的映射不单射，因此没有整体双侧逆。
\end{solution}

\begin{solution}{one-dimensional-local-global-inverse}
由中值定理，若 $x_1<x_2$，则
$f(x_2)-f(x_1)=f'(\xi)(x_2-x_1)>0$，所以 $f$ 严格递增。连续严格单调
函数把开区间映为开区间
\[
 f((a,b))=(\lim_{x\downarrow a}f(x),\lim_{x\uparrow b}f(x)),
\]
其中端点允许为无穷。它到该像区间是连续双射，逆函数连续；因 $f'$ 从不
为零，一维逆函数定理及导数公式说明逆函数属于 $C^1$。
\end{solution}

\begin{solution}{implicit-circle-charts}
圆周由 $F(x,y)=x^2+y^2-1=0$ 给出。若 $y_0>0$，邻近图像为
$y=\sqrt{1-x^2}$；若 $y_0<0$，为 $y=-\sqrt{1-x^2}$，且
$dy/dx=-x/y$。在 $x_0>0$ 或 $x_0<0$ 的左右两侧则分别写成
$x=\pm\sqrt{1-y^2}$，且 $dx/dy=-y/x$。每一点至少有一个坐标偏导不为
零，因此这些图覆盖整个圆周。
\end{solution}

\begin{solution}{implicit-equation-expansion}
令 $F(x,y)=y+x\sin y-x^2$。在原点 $F_y=1$，故可唯一解为
$y=\varphi(x)$。对
\[
 \varphi+x\sin\varphi-x^2=0
\]
求导，得
\[
 \varphi'+\sin\varphi+x\cos\varphi\,\varphi'-2x=0.
\]
代入 $x=0$ 得 $\varphi'(0)=0$。再求导并在原点代入，只有
$\varphi''(0)-2=0$ 留下，所以 $\varphi''(0)=2$。相应地
$\varphi(x)=x^2+o(x^2)$。
\end{solution}

\begin{solution}{implicit-system-two-equations}
把左端记为 $F(u,v,x,y)$。对 $(u,v)$ 的 Jacobian 在原点为
\[
 \begin{pmatrix}1&1\\1&1\end{pmatrix},
\]
不可逆，隐函数定理不适用，而且这里确实不存在邻域内的唯一 $C^1$ 解。
由第一式 $v=-u-xy$，代入第二式得到
\[
 e^u-u-1=x(1+y).
\]
左端的导数为 \(e^u-1\)，故在负半轴严格递减、在正半轴严格递增，
在 $u=0$ 有严格极小值零，并且等于 $u^2/2+O(u^3)$。当
$x(1+y)>0$ 很小时，方程在零附近有一正一负两个解；当右端为负时没有解。
所以解既不唯一，也不能在原点附近对全部 $(x,y)$ 存在。
\end{solution}

\begin{solution}{sphere-regular-level-tangent}
$F(x)=\norm x^2$ 的导数为 $DF(x)h=2\langle x,h\rangle$。在
$S^{n-1}=F^{-1}(1)$ 上 $x\ne0$，这个线性泛函满射到 $\R$，故 $1$ 是
正则值。核公式给出
\[
 T_xS^{n-1}=\{h:\langle x,h\rangle=0\}=x^\perp.
\]
\end{solution}

\begin{solution}{matrix-determinant-level-set}
在可逆矩阵 $A$ 处，
\[
 D(\det)(A)H=\det(A)\operatorname{tr}(A^{-1}H).
\]
这个公式可由
$\det(A+tH)=\det(A)\det(I+tA^{-1}H)$ 及
$\det(I+tK)=1+t\operatorname{tr}K+o(t)$ 得到。当 $\det A=1$ 时，该导数
不是零泛函（取 $H=A$ 即可），所以 $1$ 是正则值。在单位矩阵处，
\[
 T_I\det^{-1}(1)=\{H:\operatorname{tr}H=0\}.
\]
\end{solution}

\begin{solution}{graph-tangent-kernel}
参数化 $\Psi(x)=(x,\varphi(x))$ 的导数为
$D\Psi(a)h=(h,D\varphi(a)h)$，故图像的切空间是这些向量的集合。
另一方面
\[
 DF(a,\varphi(a))(h,k)=k-D\varphi(a)h.
\]
它的核恰由 $k=D\varphi(a)h$ 描述，正好等于参数化所得切空间。
\end{solution}

\begin{solution}{singular-level-crossing}
零水平集为两条直线 $y=x$ 与 $y=-x$ 的并。原点处 $DF(0,0)=0$，不满射。
假设存在题中坐标片 \(\gamma(t)=(x(t),y(t))\)。
由 \(\gamma'(0)\ne0\)，至少一个坐标导数非零。若 \(x'(0)\ne0\)，
一元逆函数定理说明缩小参数区间后，曲线像为单值图像 \(y=\psi(x)\)。
但这个像仍是零水平集的一个相对开邻域，所以它含有充分小的
\((s,s)\) 和 \((s,-s)\)（\(s\ne0\)），不可能是这样的图像。
若 \(y'(0)\ne0\)，交换坐标得到相同矛盾。这里使用了坐标片的局部
单一参数描述；一般浸入曲线可以由不同参数支相交，不能直接排除。
\end{solution}

\begin{solution}{constant-rank-projection-model}
设线性映射 $A:\R^n\to\R^m$ 的秩为 $r$。在源空间选取
$v_1,\ldots,v_r$，使 $Av_1,\ldots,Av_r$ 构成像空间的基，再以
$\ker A$ 的基补成源空间基；在目标空间以 $Av_1,\ldots,Av_r$ 开头并
补成一组基。在这两组基下，$A$ 的矩阵为
\[
 \begin{pmatrix}I_r&0\\0&0\end{pmatrix},
\]
即 $(u,v)\mapsto(u,0)$。基变换矩阵都是可逆的。
\end{solution}

\begin{solution}{constant-rank-fibers}
在常秩坐标中，$F(u,v)=(u,0)$。若所考察的纤维非空且对应值为
$(u_0,0)$，则局部纤维正是
\[
 \{(u_0,v):v\in\R^{n-r}\},
\]
经源坐标变换的逆像。它由 $n-r$ 个参数描述，是一个 $C^1$ 局部图像；
切空间是 $DF$ 的核。
\end{solution}

\begin{solution}{rank-one-map-normal-form}
固定 $y=y_0$ 时，映射
$x\mapsto e^x(\cos y_0,\sin y_0)$ 的导数从不为零；以径向坐标作为目标
第一坐标，局部写成 $x\mapsto(x,0)$（再作一维重参数化）。固定
$x=x_0$ 时，$y\mapsto e^{x_0}(\cos y,\sin y)$ 的导数也从不为零；以
圆周切向坐标为第一坐标，同样局部化为 $y\mapsto(y,0)$。两者局部都是
秩一浸入，但前者的像沿射线，后者的像沿圆周，整体性质不同。
\end{solution}

\begin{solution}{inverse-newton-fixed-linearization}
取正文所选闭球 \(\overline B(0,r)\)、\(0\le q<1\)，并设
\(\norm y<(1-q)r\)、\(x_0\in\overline B(0,r)\)。
所有迭代点均留在该球内，压缩映射定理给出唯一不动点。
不动点 $x_*$ 满足 $f(x_*)=y$。压缩估计直接给出
\[
 \norm{x_{m+1}-x_*}\le q\norm{x_m-x_*}\le q^{m+1}\norm{x_0-x_*}.
\]
还可用相邻差估计消去未知的 $x_*$：
\[
 \norm{x_m-x_*}
 \le\sum_{j=m}^{\infty}\norm{x_{j+1}-x_j}
 \le\frac{q^m}{1-q}\norm{x_1-x_0}.
\]
上面是先验误差界。若使用已经算出的相邻差，则对 \(m\ge1\)，
\[
 \norm{x_m-x_*}\le\frac{q}{1-q}\norm{x_m-x_{m-1}},
\]
因为后续相邻差分别不超过
\(q\norm{x_m-x_{m-1}},q^2\norm{x_m-x_{m-1}},\ldots\)。
这才是可以在迭代中直接检查的后验误差界。
\end{solution}

\begin{solution}{inverse-necessity-counterexamples}
连续双射但导数奇异：$f(x)=x^3$ 是 $\R$ 上连续双射，且 $f'(0)=0$；其逆
$y^{1/3}$ 在零点不可微，这也同时给出第三种现象。局部可逆但不全局单射：
$f(x,y)=(e^x\cos y,e^x\sin y)$ 的 Jacobian 处处可逆，却以 $2\pi$ 为
$y$ 的周期。这些例子分别表明，拓扑双射、局部微分同胚和整体微分同胚是
不同层次的条件。
\end{solution}

\begin{solution}{hadamard-one-dimensional-analogue}
$f'$ 连续且从不为零。由于 $\R$ 连通，$f'$ 不可能改变符号；否则介值定理
会给出零点。因此 $f$ 严格单调。条件 $|f(x)|\to\infty$ 排除两端趋向同一
方向；若 $f$ 递增，则必有 $f(x)\to-\infty$ 当 $x\to-\infty$，且
$f(x)\to\infty$ 当 $x\to\infty$，故像为 $\R$。递减情形类似。
所以 $f$ 是 $\R$ 到 $\R$ 的双射，局部逆公式拼合成全局 $C^1$ 逆。
\end{solution}
